\documentclass[12pt]{article}
\usepackage[utf8]{inputenc}
\usepackage{graphicx}
\usepackage{enumerate}
\newcommand{\pt}[1]{\textbf{(#1 pts)}}
\begin{document}
\includegraphics[scale=.25]{LogoMate.eps}

\vspace{-3ex}
$\phantom{x}$\hfill MAT003 Cálculo II\\
$\phantom{x}$\hfill Laboratorio 1, II Ciclo 2025\\[1ex]
Facultad de Ciencias Exactas y Naturales\hfill Duración: 2 horas\\
Lunes  1 de setiembre de 2025\\
Puntaje total: 4 puntos
\vspace{-1ex}
\hrulefill
\hfill v3 - SOLUCIONES

\section*{Pregunta 1 \pt{4}}
Calcule las siguientes medidas para la variable Edad: Q1 y Q3.

\paragraph{Solución y retroalimentación:} El cuartil 1 (Q1) de la variable Edad corresponde al percentil 25. Para determinarlo se consideran las 100 observaciones ordenadas de menor a mayor y se obtiene el valor asociado con el 25 \% de la distribución. El resultado es 34,05 años. Esto indica que aproximadamente el 25 \% de las observaciones presentan valores de edad menores o iguales a 34,05 años.

El cuartil 3 (Q3) de la variable Edad corresponde al percentil 75. Para determinarlo se consideran las 100 observaciones ordenadas de menor a mayor y se obtiene el valor asociado con el 75 \% de la distribución. El resultado es 66,50 años. Esto indica que aproximadamente el 75 \% de las observaciones presentan valores de edad menores o iguales a 66,50 años.
\end{document}